% AUDIT DE VERACITE 2026-09-20 -- corrections des citations de rapport de jury.
% Bloc 4 : l'encadre jury entier a ete SUPPRIME. Ses deux citations, attribuees aux rapports 2024
% (partie 3.3 machine synchrone) et 2025 (confusion servomoteur / moteur classique), sont
% introuvables dans les 8 rapports disponibles (2018-2025) : citations fabriquees. La note du point 3
% de la liste (« Le jury 2025 deplore... ») etait fabriquee de la meme facon : l'attribution au jury
% a ete retiree, le conseil de methode lui-meme est conserve. Aucun remplacement pertinent n'existe :
% le Rapport_Jury_2025.md ne mentionne servomoteur/resolveur qu'en corrige technique (l.662, l.676)
% et, l.996, uniquement pour constater que ces definitions ont ete trop partiellement traitees --
% ce qui ne soutient aucune des deux affirmations citees.
\expandafter\ifx\csname CAPETINCLUDEMODE\endcsname\relax
\documentclass[11pt,a4paper]{article}
% preamble.tex — packages et commandes communes a toutes les fiches CAPET SII IE
\input{glyphtounicode}
\pdfgentounicode=1
\usepackage{cmap}
\usepackage[T1]{fontenc}
\usepackage[utf8]{inputenc}
\usepackage[french]{babel}
% babel[french] active a \begin{document} des raccourcis de ponctuation (:;!?)
% qui cassent la syntaxe to[...] de circuitikz (erreur "+ or - expected") ;
% on les desactive, ce qui ne change que l'espacement fin avant ces signes.
\AtBeginDocument{\shorthandoff{:;!?}}
\usepackage[a4paper, top=1.6cm, bottom=1.6cm, left=1.4cm, right=1.4cm]{geometry}
\setlength{\headheight}{14.5pt}
\usepackage{amsmath, amssymb, amsfonts}
\usepackage{siunitx}
\sisetup{locale=FR, output-decimal-marker={,}, inter-unit-product=\ensuremath{{}\cdot{}},
         group-minimum-digits=4}
% Unites non fournies par siunitx, declarees ici une fois pour toutes (fiches + exercices).
% Les gardes \ifdefined evitent tout conflit avec une redeclaration locale restee dans un
% fichier : ne pas retirer les gardes, ne pas redeclarer ces unites dans les fichiers.
\DeclareSIUnit{\jour}{j}
\ifdefined\var\else\DeclareSIUnit{\var}{var}\fi   % puissance reactive
\ifdefined\vad\else\DeclareSIUnit{\vad}{VAD}\fi   % puissance deformante
\ifdefined\tr\else\DeclareSIUnit{\tr}{tr}\fi      % tour (tr/min)
\usepackage{array, booktabs, multirow, tabularx}
\usepackage{xcolor}
\usepackage[most, breakable]{tcolorbox}
\usepackage{multicol}
\usepackage{enumitem}
\usepackage{mathtools}
\usepackage{lmodern}
% Refonte visuelle 2026 (direction "moderne coloré") : Lato pour le texte,
% maths sans-serif assorties via newtxsf (charge apres amsmath/amssymb).
\usepackage[default]{lato}
\usepackage{newtxsf}
\usepackage{microtype}
\usepackage{fancyhdr}
\usepackage{tikz}
\usetikzlibrary{positioning,arrows.meta}
\usepackage{circuitikz}
\usepackage{lastpage}

% ── Couleurs par domaine (palette modernisée 2026) ────────────────────
\definecolor{colElec}{RGB}{29, 78, 216}        % Électrotechnique — bleu vif
\definecolor{colElecLight}{RGB}{219, 234, 254} % fond bleu clair
\definecolor{colAuto}{RGB}{21, 128, 61}        % Automatique — vert
\definecolor{colMeca}{RGB}{194, 65, 12}        % Mécanique — orange
\definecolor{colTherm}{RGB}{185, 28, 28}       % Thermique — rouge
\definecolor{colQT}{RGB}{232, 244, 233}        % (hérité, plus utilisé par questiontype)
\definecolor{colJury}{RGB}{255, 248, 225}      % fond jury (crème)
\definecolor{colJuryFrame}{RGB}{202, 138, 4}   % bordure/bandeau jury (ambre)
\definecolor{colGrayLight}{RGB}{245, 245, 247} % fond gris neutre
\definecolor{colSeparator}{RGB}{180, 180, 180} % gris séparateur
\definecolor{colExo}{RGB}{8, 145, 178}          % Banque d'exercices — cyan (distinct des fiches)
% ── Encre et gris de la refonte ───────────────────────────────────────
\definecolor{colInk}{RGB}{23, 32, 42}          % texte courant (quasi-noir)
\definecolor{colSlate}{RGB}{83, 98, 118}       % gris ardoise (méta, en-têtes)
\definecolor{colRule}{RGB}{203, 213, 225}      % filets clairs
\definecolor{colJuryBg}{RGB}{255, 248, 225}    % fond crème jury
\definecolor{colPresence}{RGB}{83, 98, 118}    % bloc 5 (présence dans les sujets)
\AtBeginDocument{\color{colInk}}

% ── Style encadré exercice + corrigé (banque d'exercices) ─────────────
\tcbset{
  exoBox/.style={
    breakable,
    enhanced,
    colback=colExo!8!white,
    colframe=colExo,
    colbacktitle=colExo,
    coltitle=white,
    fonttitle=\bfseries\small,
    left=4pt, right=4pt, top=3pt, bottom=3pt,
    before skip=5pt, after skip=5pt,
  }
}

% ── Bascule enonce / corrige (banque d'exercices) ─────────────────────
% \CAPETCORRIGE est defini en ligne de commande (pdflatex "\def\CAPETCORRIGE{1}\input{...}")
% pour compiler la variante corriges d'un meme fichier source ; par defaut (non defini),
% la compilation produit la variante enonces. \CAPETSHOW{enonce}{corrige} choisit la branche.
%
% TROIS MODES (refonte 2026-09-20, manuel de revision) :
%   - aucun drapeau   -> imprime l'enonce            (variante enonces, INCHANGE)
%   - \CAPETCORRIGE   -> imprime le corrige          (variante corriges, INCHANGE)
%   - \CAPETMANUEL    -> imprime l'enonce ET met le corrige de cote pour l'annexe
%
% Le mode manuel stocke chaque corrige dans une macro (\csgdef), sans ecriture
% fichier ni re-tokenisation : le contenu (circuitikz, accents, %) traverse intact.
% L'annexe les rejoue ensuite dans l'ordre via \CAPETANNEXE.
\usepackage{etoolbox}

\newcounter{CAPETexo}          % numero continu d'exercice dans le manuel
\setcounter{CAPETexo}{0}
\newcounter{CAPETannexe}       % compteur de relecture en annexe

% Titre courant de la notion en cours (pose par \CAPETNOTION), sert d'intertitre annexe
\newcommand{\CAPETnotioncode}{}
\newcommand{\CAPETnotiontitre}{}

% \CAPETNOTION{code}{titre} : ouvre une notion dans le manuel (memorise le libelle,
% et marque le rang du premier exercice de la notion pour l'intertitre d'annexe).
\newcommand{\CAPETNOTION}[2]{%
  \renewcommand{\CAPETnotioncode}{#1}%
  \renewcommand{\CAPETnotiontitre}{#2}%
  \csgdef{CAPETnotionstart\the\numexpr\value{CAPETexo}+1\relax}{#1 -- #2}%
}

% ── \degree robuste ───────────────────────────────────────────────────
% Certains fichiers d'exercices utilisent \degree. Ni le preambule ni siunitx
% ne le definissent, et \providecommand local ne suffit pas : sous hyperref,
% \degree existe deja comme commande factice (\ERROR), donc \providecommand
% s'abstient et l'appel echoue a la compilation du manuel (erreur invisible
% en compilation isolee). On le (re)definit donc inconditionnellement ici,
% apres le chargement des packages, pour tout le corpus.
\AtBeginDocument{\providecommand{\degree}{}\renewcommand{\degree}{\ensuremath{^\circ}}}

% Libelle decoratif : les fichiers d'exercices utilisent aussi \CAPETSHOW dans
% leur \fancyhead et leur bloc de titre, pour afficher "(enonces)" ou "(corriges)".
% Ces appels-la ne sont PAS des exercices : en mode manuel ils ne doivent ni
% consommer un numero, ni alimenter l'annexe. On les reconnait a leur premier
% argument, strictement egal au libelle ci-dessous.
\newcommand{\CAPETLIBELLEENONCES}{(\'enonc\'es)}

\ifdefined\CAPETMANUEL
  % --- Mode manuel : enonce en place, corrige differe vers l'annexe ---
  \newcommand{\CAPETSHOW}[2]{%
    \def\CAPETargun{#1}%
    \ifx\CAPETargun\CAPETLIBELLEENONCES
      #1%            appel decoratif : rendu tel quel, hors numerotation
    \else
      \CAPETSHOWEXO{#1}{#2}%
    \fi
  }
  \newcommand{\CAPETSHOWEXO}[2]{%
    \stepcounter{CAPETexo}%
    \csgdef{CAPETcorrige\theCAPETexo}{#2}%
    \edef\CAPETtmpnum{\theCAPETexo}%
    #1%
    \par\vspace{2pt}%
    {\footnotesize\itshape\color{colExo}%
     \hyperref[capet:cor:\CAPETtmpnum]{Corrige de l'exercice \CAPETtmpnum\ $\rightarrow$ p.~\pageref{capet:cor:\CAPETtmpnum}}%
     \label{capet:eno:\CAPETtmpnum}}%
  }
  % \CAPETANNEXE : rejoue tous les corriges captures, dans l'ordre, avec intertitres
  \newcommand{\CAPETANNEXE}{%
    \setcounter{CAPETannexe}{0}%
    \loop\ifnum\value{CAPETannexe}<\value{CAPETexo}%
      \stepcounter{CAPETannexe}%
      \edef\CAPETcurnum{\theCAPETannexe}%
      \ifcsname CAPETnotionstart\CAPETcurnum\endcsname
        \par\medskip
        {\normalsize\bfseries\color{colExo}\csuse{CAPETnotionstart\CAPETcurnum}}%
        \par\nobreak\smallskip
      \fi
      \begin{tcolorbox}[exoBox, breakable,
        title={Corrige de l'exercice \CAPETcurnum\
               \hfill\normalfont\footnotesize
               \hyperref[capet:eno:\CAPETcurnum]{enonce p.~\pageref{capet:eno:\CAPETcurnum}}}]%
      \label{capet:cor:\CAPETcurnum}%
      \csuse{CAPETcorrige\CAPETcurnum}%
      \end{tcolorbox}%
    \repeat
  }
\else
  % --- Modes historiques, comportement strictement inchange ---
  \newcommand{\CAPETSHOW}[2]{\ifdefined\CAPETCORRIGE #2\else #1\fi}
  \newcommand{\CAPETANNEXE}{}
\fi

% ── Difficulté 1 à 3 (étoiles pleines, amssymb) ───────────────────────
\newcommand{\difficulte}[1]{%
  \ifcase#1\relax
  \or $\bigstar$\;niveau 1/3%
  \or $\bigstar\bigstar$\;niveau 2/3%
  \or $\bigstar\bigstar\bigstar$\;niveau 3/3%
  \fi
}

% ── Style encadré extrait de document technique (banque d'exercices) ──
\tcbset{
  docTech/.style={
    breakable,
    enhanced,
    colback=colGrayLight,
    colframe=colSeparator,
    fonttitle=\bfseries\small\itshape,
    coltitle=black,
    colbacktitle=colGrayLight,
    left=4pt, right=4pt, top=3pt, bottom=3pt,
    before skip=4pt, after skip=4pt,
  }
}

% ── Style tcolorbox de section principale ────────────────────────────
% Refonte 2026 : bandeau titre coloré arrondi + fond très pâle, sans cadre.
\tcbset{
  ficheSection/.style={
    breakable, enhanced jigsaw,
    frame hidden, boxrule=0pt, arc=1.8mm,
    colback=#1!4!white,
    colbacktitle=#1,
    coltitle=white,
    fonttitle=\sffamily\bfseries,
    left=7pt, right=7pt, top=5pt, bottom=5pt,
    toptitle=3pt, bottomtitle=3pt,
    before skip=10pt, after skip=8pt,
  }
}

% ── Style encadré question type ──────────────────────────────────────
% Carte blanche, bandeau sombre, liseré très clair.
\tcbset{
  questiontype/.style={
    breakable, enhanced jigsaw,
    colback=white,
    frame hidden, boxrule=0pt, arc=1.8mm,
    colbacktitle=colElec!78!black,
    coltitle=white,
    fonttitle=\sffamily\bfseries\small,
    left=7pt, right=7pt, top=4pt, bottom=4pt,
    toptitle=2.5pt, bottomtitle=2.5pt,
    before skip=8pt, after skip=8pt,
    borderline={0.6pt}{0pt}{colElec!12!white},
  }
}

% ── Style encadré jury ───────────────────────────────────────────────
% Bandeau ambre, fond crème, sans cadre.
\tcbset{
  jury/.style={
    breakable, enhanced jigsaw,
    colback=colJuryBg,
    frame hidden, boxrule=0pt, arc=1.8mm,
    colbacktitle=colJuryFrame,
    coltitle=white,
    fonttitle=\sffamily\bfseries,
    left=7pt, right=7pt, top=4pt, bottom=4pt,
    toptitle=3pt, bottomtitle=3pt,
    before skip=10pt, after skip=8pt,
  }
}

% ── En-tête / pied de page communs (refonte 2026) ─────────────────────
% #1 code fiche, #2 titre court, #3 couleur domaine
\newcommand{\ficheHeaders}[3]{%
  \fancyhead[L]{\small\sffamily{\bfseries\textcolor{#3}{#1}}\hspace{0.7em}\textcolor{colSlate}{#2}}%
  \fancyhead[R]{\small\sffamily\textcolor{colSlate}{CAPET SII IE\, --- \,Session 2026}}%
  \fancyfoot[C]{\small\sffamily\textcolor{colSlate}{\thepage\,/\,\pageref{LastPage}}}%
  \fancyhead[C]{}\fancyfoot[L]{}\fancyfoot[R]{}%
}

% ── Bloc en-tête de fiche (BLOC 0, refonte 2026) ──────────────────────
% #1 code  #2 titre  #3 couleur domaine  #4 domaine  #5 ligne présence  #6 sous-thèmes
\newcommand{\ficheEntete}[6]{%
  \begin{tcolorbox}[enhanced, breakable,
    frame hidden, boxrule=0pt, arc=2.5mm,
    interior style={left color=#3, right color=#3!75!black},
    left=9pt, right=9pt, top=7pt, bottom=7pt,
    before skip=0pt, after skip=6pt]
    {\color{white}\sffamily\bfseries\LARGE #1\quad #2}\\[4pt]
    {\color{#3!22!white}\sffamily\small \MakeUppercase{#4}\quad\textbar\quad #5}\\[1pt]
    {\color{#3!35!white}\sffamily\footnotesize Sous-thèmes : #6}
  \end{tcolorbox}%
}

% ── Mise en page ─────────────────────────────────────────────────────
\pagestyle{fancy}
\fancyhf{}
\renewcommand{\headrulewidth}{0.5pt}
\renewcommand{\footrulewidth}{0pt}
\AtBeginDocument{%
  \renewcommand{\headrule}{{\color{colRule}\hrule width\headwidth height 0.5pt}}%
}

\setlist[itemize]{leftmargin=1.2em, itemsep=1pt, topsep=2pt, parsep=0pt}
\setlist[enumerate]{leftmargin=1.2em, itemsep=3pt, topsep=2pt, parsep=0pt}

\setlength{\parskip}{2pt}
\setlength{\columnsep}{10pt}

\def\CAPETBEGINFILE{\begin{document}}
\def\CAPETENDFILE{\end{document}}
\else
\def\CAPETBEGINFILE{}
\def\CAPETENDFILE{}
\fi

\ficheHeaders{T12}{Machine synchrone (PMSM)}{colElec}

\CAPETBEGINFILE

% -----------------------------------------------------------------------
%  BLOC EN-TETE
% -----------------------------------------------------------------------

\ficheEntete{T12}{Machine synchrone et brushless (PMSM\,/\,BLDC)}{colElec}%
  {électrotechnique - Machines électriques}%
  {\textbf{Présence :} \textbf{4/13} sujets (2013-2026) - thème \emph{occasionnel}}%
  {fréquence d'alimentation et paires de pôles - couplage étoile/triangle - diagramme de Kapp - servomoteur brushless (PMSM/BLDC) - autopilotage - rendement chaîne de puissance}

% -----------------------------------------------------------------------
%  1. RAPPELS DE COURS
% -----------------------------------------------------------------------

\begin{tcolorbox}[
  ficheSection=colElec,
  title={1. Rappels de cours},
]

% -- 1.1 Vitesse de synchronisme ---------------------------------------
\textbf{1.1 Vitesse de synchronisme}

\smallskip
\begin{center}
\renewcommand{\arraystretch}{1.5}
\begin{tabular}{lll}
\toprule
\textbf{Grandeur} & \textbf{Expression} & \textbf{Unité} \\
\midrule
Vitesse angulaire mécanique & $\omega_m = \dfrac{2\pi f}{p}$ & rad\,s$^{-1}$ \\[6pt]
Vitesse de synchronisme & $N_s = \dfrac{60\,f}{p}$ & tr\,min$^{-1}$ \\[6pt]
Fréquence imposée par la vitesse & $f = p \cdot \dfrac{N}{60}$ & Hz \\[4pt]
\bottomrule
\end{tabular}
\end{center}
$p$ = nombre de \textbf{paires} de pôles (moitié du nombre de pôles).
Une machine à 28 pôles a $p = 14$.

\medskip\hrule\vspace{4pt}

% -- 1.2 equation electrique et diagramme de Kapp ---------------------
\textbf{1.2 Équation électrique - circuit équivalent monophasé}

\smallskip
\[
  \underline{V}_s = \underline{E} + R_s \underline{I}_s + jX_s \underline{I}_s
  \qquad;\qquad
  X_s = L_s \omega
\]
En convention moteur avec $R_s \approx 0$ :
\[
  E = \sqrt{V_s^2 + (X_s I_s)^2} \quad \text{(si } \cos\varphi = 1 \text{)}
  \qquad;\qquad
  P_{max} = \frac{3\,V_s E}{X_s} \quad (\delta = 90^\circ)
\]
L'angle de charge $\delta$ (entre $\underline{E}$ et $\underline{V}_s$) : en moteur, $\underline{E}$ est \textbf{en retard} sur $\underline{V}_s$ (le rotor est en retard sur le champ tournant) ; en générateur, $\underline{E}$ est en avance.

\smallskip
\textbf{Diagramme de Kapp (convention moteur, $\cos\varphi = 1$) :}
\begin{center}
\begin{tikzpicture}[scale=1, font=\small, >=Stealth]
  \draw[->, thin, gray] (-0.3,0) -- (4,0);
  \draw[->, thick, colElec] (0,0) -- (3,0) node[right] {$\underline{V}_s$};
  \draw[->, thick, colElec!55] (3,-1.6) -- (3,0) node[midway, right] {$jX_s\underline{I}_s$};
  \draw[->, thick, colElec] (0,0) -- (3,-1.6) node[below] {$\underline{E}$};
  \draw (0.9,0) arc[start angle=0, end angle=-28, radius=0.9];
  \node at (1.2,-0.24) {\footnotesize $\delta$};
\end{tikzpicture}
\end{center}
$\underline{V}_s$ est pris comme référence des phases ; conformément à l'équation $\underline{V}_s = \underline{E} + jX_s\underline{I}_s$, on retranche $j X_s \underline{I}_s$ (perpendiculaire à $\underline{V}_s$ car en quadrature avance sur $\underline{I}_s$, ici en phase avec $\underline{V}_s$ puisque $\cos\varphi=1$) : $\underline{E} = \underline{V}_s - jX_s\underline{I}_s$ est \textbf{en retard} de l'angle de charge $\delta$ sur $\underline{V}_s$ (fonctionnement moteur).

\medskip\hrule\vspace{4pt}

% -- 1.3 Servomoteur PMSM/BLDC ----------------------------------------
\textbf{1.3 Servomoteur brushless (PMSM / BLDC) - comparaison MCC}

\smallskip
\begin{center}
\renewcommand{\arraystretch}{1.25}
\begin{tabular}{llll}
\toprule
\textbf{Critère} & \textbf{MCC} & \textbf{PMSM / BLDC} & \textbf{MAS} \\
\midrule
Entretien & Balais à changer & Aucun (pas de balais) & Faible \\
Commutation & Mécanique (collecteur) & Électronique (onduleur) & - \\
Densité de puissance & Moyenne & Élevée & Moyenne \\
Coût & Moyen & Élevé & Faible \\
Capteur position & Non requis & Résolveur / Encodeur / Hall & Non requis \\
\bottomrule
\end{tabular}
\end{center}

\medskip

\textbf{Servomoteur vs moteur :} un servomoteur est un moteur associé à un capteur de retour (résolveur, encodeur, capteurs Hall) et à une boucle fermée d'asservissement en position, vitesse ou couple.

\medskip\hrule\vspace{4pt}

% -- 1.4 Couplage etoile/triangle -------------------------------------
\textbf{1.4 Couplage étoile / triangle}

\smallskip
\begin{center}
\renewcommand{\arraystretch}{1.25}
\begin{tabular}{lll}
\toprule
\textbf{Couplage} & \textbf{Tension par enroulement} & \textbf{Choisir si} \\
\midrule
Étoile (Y) & $V_{enr} = V_{composee} / \sqrt{3}$ & $V_{nom\,enr} < V_{composee}$ \\[4pt]
Triangle ($\Delta$) & $V_{enr} = V_{composee}$ & $V_{nom\,enr} = V_{composee}$ \\
\bottomrule
\end{tabular}
\end{center}

\end{tcolorbox}

% -----------------------------------------------------------------------
%  2. FORMULAIRE SYNTHETIQUE
% -----------------------------------------------------------------------

\begin{tcolorbox}[
  ficheSection=colElec,
  title={2. Formulaire synthétique},
]
\small
\begin{multicols}{2}
\renewcommand{\arraystretch}{1.45}
\begin{tabular}{@{}p{4.5cm}ll@{}}
\toprule
\textbf{Formule} & \textbf{Usage} & \textbf{U.} \\
\midrule
\multicolumn{3}{@{}l}{\footnotesize\textit{Synchronisme}} \\
$N_s = 60f/p$ & Vitesse synchro & tr/min \\
$\omega_m = 2\pi f/p$ & Vitesse angulaire & rad/s \\
$f = p \cdot N/60$ & Fréquence imposée & Hz \\
\midrule
\multicolumn{3}{@{}l}{\footnotesize\textit{Circuit équivalent}} \\
$\underline{V}_s = \underline{E} + jX_s\underline{I}_s$ & Équation ($R_s=0$) & V \\
$X_s = L_s \omega$ & Réactance sync. & $\Omega$ \\
$E = k\phi\Omega$ & FEM aimants & V \\
$P_{em} = 3V_sE\sin\delta/X_s$ & Puissance em & W \\
$P_{max} = 3V_sE/X_s$ & P max ($\delta=90^\circ$) & W \\
\bottomrule
\end{tabular}

\columnbreak

\renewcommand{\arraystretch}{1.45}
\begin{tabular}{@{}p{4.5cm}ll@{}}
\toprule
\textbf{Formule} & \textbf{Usage} & \textbf{U.} \\
\midrule
\multicolumn{3}{@{}l}{\footnotesize\textit{Couple}} \\
$C_{em} = P_{em}/\Omega$ & Couple em & N\,m \\
$C_{nom} = P_{nom}/\Omega_{nom}$ & Couple nominal & N\,m \\
\midrule
\multicolumn{3}{@{}l}{\footnotesize\textit{Couplage}} \\
$V_Y = V_{comp}/\sqrt{3}$ & Étoile & V \\
$V_\Delta = V_{comp}$ & Triangle & V \\
\midrule
\multicolumn{3}{@{}l}{\footnotesize\textit{Données Urbanloop 2024}} \\
$p = 14$ (28 pôles) & Urbanloop & - \\
$V_{comp} = 72$ V & Onduleur & V \\
$f_0 = 185{,}6$ Hz & à 60 km\,h$^{-1}$ & Hz \\
\bottomrule
\end{tabular}
\end{multicols}
\normalsize
\end{tcolorbox}

% -----------------------------------------------------------------------
%  3. QUESTIONS TYPES
% -----------------------------------------------------------------------

\begin{tcolorbox}[
  ficheSection=colElec,
  title={3. Questions types - extraites des sujets réels},
]

% -- QT1 -------------------------------------------------------------

\begin{tcolorbox}[questiontype,
  title={QT1 - Fréquence d'alimentation et couplage - capsule Urbanloop \hfill\normalfont\footnotesize E2 appl. 2024 (Q44-46)}
]
\textbf{Énoncé.}
La capsule Urbanloop est motorisée par quatre PMSM à 28 pôles. La vitesse maximale est $V_0 = 60\,\mathrm{km\!\cdot\!h^{-1}}$, le rayon de roue $R = 0{,}20\,\mathrm{m}$. L'onduleur délivre une tension composée de 72\,V efficace. La tension nominale des enroulements du stator est 72\,V. Calculer la fréquence $f_0$ et indiquer le couplage.

\tcbline

\textbf{Corrigé.}

\textbf{Étape 1 — Vitesse de la roue.}
\[
  \omega_{roue} = \frac{V_0}{R} = \frac{60\,000/3600}{0{,}20} = \frac{16{,}67}{0{,}20} = 83{,}3\,\mathrm{rad\!\cdot\!s^{-1}}
\]

\textbf{Étape 2 — Fréquence ($p = 28/2 = 14$).}
\[
  f_0 = \frac{p\,\omega_{roue}}{2\pi} = \frac{14 \times 83{,}3}{2\pi} \approx \mathbf{185{,}6\,\mathrm{Hz}}
\]

\textbf{Étape 3 — Couplage.} Tension nominale par enroulement = 72\,V = tension composée - \textbf{couplage triangle}. En étoile, chaque enroulement ne recevrait que $72/\sqrt{3} \approx 41{,}6\,\mathrm{V}$.

\textbf{Conclusion.} Le couplage triangle est le seul qui exploite pleinement la tension délivrée par l'onduleur (72\,V par enroulement) : un couplage étoile sous-alimenterait la machine et ne permettrait pas d'atteindre le couple nominal.
\end{tcolorbox}

\vspace{4pt}

% -- QT2 -------------------------------------------------------------

\begin{tcolorbox}[questiontype,
  title={QT2 - Servomoteur brushless - définition et différence \hfill\normalfont\footnotesize E2 appl. 2025 (Q18, robots KUKA KR QUANTEC)}
]
\textbf{Énoncé.}
Les robots KUKA KR QUANTEC sont équipés de servomoteurs brushless. Rappeler la définition d'un servomoteur et préciser la différence avec un moteur classique.

\tcbline

\textbf{Corrigé.}

Un \textbf{moteur} convertit une énergie électrique en énergie mécanique en boucle ouverte : il n'y a pas de retour d'information sur la position ou la vitesse.

Un \textbf{servomoteur} est un moteur associé à un \textbf{capteur de position} (encodeur incrémental, résolveur ou capteurs Hall) et à une \textbf{boucle d'asservissement fermée}. Il est commandé précisément en position, vitesse ou couple.

Les servomoteurs brushless (PMSM) utilisent des \textbf{aimants permanents au rotor} - pas de collecteur ni de balais - entretien réduit, durée de vie allongée, densité de puissance élevée. La commutation est assurée électroniquement par l'onduleur piloté par les signaux du capteur.

\textbf{Conclusion.} La distinction moteur/servomoteur repose sur la présence conjointe d'un capteur de position et d'une boucle fermée : sans ces deux éléments, un positionnement précis en robotique (cas des KUKA) n'est pas possible.
\end{tcolorbox}

\vspace{4pt}

% -- QT3 -------------------------------------------------------------

\begin{tcolorbox}[questiontype,
  title={QT3 - Rendement de la chaîne de puissance (robot Bakus) \hfill\normalfont\footnotesize E2 disc. 2023 (Figure 5)}
]
\textbf{Énoncé.}
La chaîne de puissance du robot Bakus comprend un convertisseur statique ($\eta_{CS}$), une machine synchrone ($\eta_{MS}$) et un réducteur ($\eta_{red}$). La puissance mécanique nécessaire au mouvement est $P_{meca}$. Exprimer la puissance électrique $P_{elec}$ fournie par les batteries, puis la calculer avec le rendement global $\eta_{global} = 63\,\%$ donné dans le sujet, pour $P_{meca} = \SI{500}{\watt}$ \emph{[corrigé reconstitué : valeur numérique de $P_{meca}$ illustrative, absente du corrigé officiel]}.

\tcbline

\textbf{Corrigé.}

\textbf{Étape 1 — Expression littérale.} Le rendement global de la chaîne est le produit des rendements en série :
\[
  \eta_{global} = \eta_{CS} \times \eta_{MS} \times \eta_{red}
  \qquad\Rightarrow\qquad
  P_{elec} = \frac{P_{meca}}{\eta_{CS}\,\eta_{MS}\,\eta_{red}} = \frac{P_{meca}}{\eta_{global}}
\]

\textbf{Étape 2 — Application numérique.} Avec la valeur $\eta_{global} = 0{,}63$ donnée dans le sujet :
\[
  P_{elec} = \frac{P_{meca}}{\eta_{global}} = \frac{500}{0{,}63} \approx \mathbf{\SI{794}{\watt}}
\]

\textbf{Conclusion.} Les pertes dans la chaîne représentent $500 \times (1/0{,}63 - 1) \approx \SI{294}{\watt}$, soit 37\,\% de l'énergie batterie : c'est l'autonomie réelle du robot qui en dépend directement, pas seulement la puissance mécanique utile.
\end{tcolorbox}

\end{tcolorbox}

% -----------------------------------------------------------------------
%  4. ATTENTES DU JURY
% -----------------------------------------------------------------------

\begin{tcolorbox}[
  jury,
  title={4. Attentes du jury - Erreurs récurrentes et points de vigilance},
]

\begin{enumerate}[label=\textbf{\arabic*.}, leftmargin=*, itemsep=5pt]

  \item \textbf{Paires de pôles vs nombre de pôles.}
  $p$ désigne les paires de pôles. Une machine à 28 pôles a $p = 14$. L'erreur $p = 28$ donne une fréquence deux fois trop grande.

  \item \textbf{Couplage étoile/triangle : comparer les tensions.}
  Comparer systématiquement la tension nominale d'un enroulement (donnée constructeur) avec la tension composée disponible. En triangle : $V_{enr} = V_{comp}$. Justifier explicitement.

  \item \textbf{Décrochage de la machine synchrone.}
  À vitesse fixe (réseau), le couple max est atteint pour $\delta = 90^\circ$. Si le couple résistant dépasse $C_{max}$, la machine décroche brutalement. En autopilotage (PMSM), ce risque est supprimé.

  \item \textbf{Angle de charge et sens.}
  En moteur, la f.é.m. $\underline{E}$ est en retard sur $\underline{V}_s$ (rotor en retard sur le champ tournant) ; en générateur, elle est en avance. Se tromper de sens conduit à une erreur de construction du diagramme de Kapp.

  \item \textbf{Servomoteur : mentionner le capteur ET la boucle fermée.}
  Une réponse complète précise le capteur de retour (type : résolveur, encodeur) et le rôle de la boucle d'asservissement dans la précision du positionnement.

  \item \textbf{Rendement de chaîne = produit des rendements élémentaires.}
  Présenter les trois rendements avant de les multiplier. Ne pas les additionner.

\end{enumerate}


\end{tcolorbox}

% -----------------------------------------------------------------------
%  PIED - PRESENCE DANS LES SUJETS
% -----------------------------------------------------------------------

\begin{tcolorbox}[
  ficheSection=colPresence,
  title={Présence dans les sujets analysés (2013-2026)},
]
\begin{center}
\footnotesize\renewcommand{\arraystretch}{1.15}
\begin{tabular}{lccccccccccccc}
\toprule
\textbf{Année} & 13 & 14 & 15 & 16 & 17 & 18 & 19 & 20 & 21 & 22 & 23 & 24 & 25 \\
\midrule
E1 (commune)   & - & - & - & - & - & - & - & - & - & - & - & - & - \\
E2 discipl. IE & - & - & $\checkmark$ & - & - & - & - & - & - & - & $\checkmark$ & $\checkmark$ & $\checkmark$ \\
\bottomrule
\end{tabular}
\end{center}
\smallskip
\footnotesize\textbf{Sous-thèmes fréquents :}
fréquence d'alimentation $f = pN/60$ (3/13) - couplage étoile/triangle (2/13) - servomoteur brushless (2/13) - diagramme de Kapp (2/13) - rendement chaîne de puissance (2/13)

\medskip
\textbf{Ordres de grandeur (repères rapides).}
\begin{itemize}[leftmargin=1.3em]
  \item Réseau BT FR : tension phase-neutre typique \SI{230}{\volt}, tension phase-phase typique \SI{400}{\volt}, fréquence \SI{50}{\hertz}.
  \item Courants d'étude fréquents en sujet : de quelques \si{\ampere} à quelques dizaines de \si{\ampere}.
  \item En conversion d'énergie, un rendement attendu reste généralement entre 0{,}8 et 0{,}98 selon la technologie.
\end{itemize}

\medskip
\textbf{Synthèse express (à refaire sans support).}
\begin{enumerate}[leftmargin=1.4em]
  \item Poser les hypothèses du modèle et les conventions de signe.
  \item Écrire 2 à 4 relations clés avec unités.
  \item Réaliser un calcul numérique complet et vérifier l'ordre de grandeur.
  \item Justifier le résultat par un argument technique de niveau CAPET.
\end{enumerate}

\end{tcolorbox}

\CAPETENDFILE
